% foundations/glossary.tex

\chapter*{Glossary}
\addcontentsline{toc}{chapter}{Glossary}

\section*{Core Concepts}

\begin{description}

\item[Alignment]
A structured representation of railway track geometry defined primarily
through curvature \(\kappa(s)\), together with associated state
variables such as position, orientation, and cant.
An alignment is not merely a curve, but a system linking geometry,
dynamics, and realization.

\item[Curvature \(\kappa(s)\)]
The fundamental geometric quantity describing the rate of change of
direction along the track.
Within the AIM framework, curvature is the primary design variable.

\item[Cant \(\phi(s)\)]
The cross-level or superelevation of the track, represented as rotation
of the track plane around its longitudinal axis.
Cant is dynamically coupled to curvature via equilibrium conditions.

\item[State]
The physical description of the track at arc-length position \(s\),
typically including position \(\gamma(s)\), orientation \(\theta(s)\),
curvature \(\kappa(s)\), and cant \(\phi(s)\).

\end{description}

\section*{Modeling and Representation}

\begin{description}

\item[Parametric Model]
A low-dimensional representation of curvature or cant defined by a small
set of parameters \(\mathbf{T}\), such as transition types or segment
lengths.

\item[Sample-Based Model]
A discrete representation of curvature or geometry defined at a set of
points \(\{s_i\}\), typically used for measurement data or numerical
optimization.

\item[Hybrid Model]
A combination of parametric structure and discrete corrections:
\[
\kappa(s)
=
\kappa_{\mathrm{param}}(s;\mathbf{T})
+
\delta\kappa(s).
\]
It allows structured modelling with local flexibility.

\item[Sparse Representation]
A minimal representation of an alignment using essential parameters,
such as curvature segments and transitions, enabling efficient storage
and computation.

\end{description}

\section*{Operators}

\begin{description}

\item[Sampling Operator \(\mathcal{S}\)]
Maps a continuous function to discrete values:
\[
\mathcal{S}(\kappa)=\{\kappa(s_i)\}.
\]

\item[Interpolation Operator \(\mathcal{I}\)]
Reconstructs a continuous curve or function from discrete data.

\item[Mechanical Operator \(\mathcal{M}\)]
Represents the physical response of the track structure, including
smoothing and deformation effects.

\item[Realization Operator \(\mathcal{R}\)]
Maps a theoretical alignment to its physically realized form:
\[
\mathcal{R}
=
\mathcal{M}\circ\mathcal{I}\circ\mathcal{S}.
\]

\item[Adjustment Operator \(\mathcal{T}\)]
A local operator representing one construction or maintenance step,
for example tamping:
\[
\gamma_{k+1}
=
\mathcal{T}(\gamma_k).
\]

\end{description}

\section*{Spaces and Transformations}

\begin{description}

\item[Track Space]
A coordinate system aligned with the track, typically parameterized by
arc-length \(s\), where geometry and dynamics are defined.

\item[World Space]
A global coordinate system, often CRS-based, used for georeferencing and
data integration.

\item[World-to-Track Transformation]
The mapping from world coordinates to intrinsic alignment coordinates,
typically requiring projection onto the alignment and evaluation of
\(s\), \(d\), \(\kappa\), and \(\phi\).

\item[Kilometre Line]
A longitudinal reference structure used for stationing, independent of
the physical geometric alignment.

\end{description}

\section*{Optimization}

\begin{description}

\item[Residual]
A component of the objective function representing deviation from a
desired property, such as smoothness, dynamic equilibrium, or data fit.

\item[Jacobian]
The matrix of derivatives of residuals with respect to parameters.
Within AIM, Jacobians are derived analytically from arc-length-based
relations.

\item[Gauss--Newton Method]
An optimization method using the approximation
\[
H \approx J^T J
\]
for solving nonlinear least-squares problems.

\item[Sequential Quadratic Programming (SQP)]
An iterative optimization framework solving a sequence of quadratic
subproblems based on linearized constraints.

\end{description}

\section*{Realization and Construction}

\begin{description}

\item[Realized Alignment]
The physically constructed track geometry, which differs from the
theoretical alignment due to discretization, construction, and mechanical
effects.

\item[Control Points]
Discrete positions \(\{s_i\}\) at which geometry is specified or
adjusted during construction or maintenance.

\item[Realization Error]
The deviation between target and realized curvature or geometry.

\item[Inverse Realization Problem]
The problem of finding a target alignment such that
\[
\mathcal{R}(\kappa_{\mathrm{target}})
\approx
\kappa_{\mathrm{desired}}.
\]

\end{description}

\section*{Advanced Geometry}

\begin{description}

\item[Geodesic Curvature \(\kappa_g\)]
The intrinsic curvature of a curve on a surface, measuring the change of
direction within the surface.
For railway alignment on the Earth, this is the physically relevant
curvature.

\item[Normal Curvature \(\kappa_n\)]
The component of curvature induced by the curvature of the underlying
surface, for example the Earth ellipsoid.

\item[Ellipsoidal Alignment]
An alignment defined directly on a reference ellipsoid rather than in a
projected coordinate system.

\item[Surface Frame]
A local coordinate frame on the ellipsoid consisting of tangent, lateral,
and surface-normal directions.

\item[Cant Rotation]
Rotation of the track cross-section around the tangent direction,
defining superelevation.

\item[Effective Lateral Acceleration]
The physically relevant lateral acceleration acting on a vehicle:
\[
a_{\mathrm{eff}}
=
v^2\kappa_g
-
\mathbf{g}\cdot\mathbf{b}_{\phi}.
\]

\end{description}

\section*{Pipeline}

\begin{description}

\item[Pipeline]
The sequence of transformations from raw data to optimized alignment:
\[
\text{data}
\rightarrow
\text{model}
\rightarrow
\text{state}
\rightarrow
\text{realization}
\rightarrow
\text{optimization}.
\]

\item[Failure Modes]
Typical breakdown points within the pipeline, such as inconsistent data,
invalid coordinate systems, unstable projection, or failed optimization.

\item[Design Principles]
Guiding principles ensuring robustness, interpretability, and consistency
across the pipeline.

\end{description}
